\section{Tutorial: Dos Dados a uma Lista de Alvos}

Este tutorial acompanha uma análise do começo ao fim: da tabela
bruta até uma lista curta de alvos de diversificação para uma região, e
depois o teste de se o método que gerou a lista prevê para onde as
regiões de fato foram. Cada passo mostra o código, o que ele imprime e
como ler o resultado. Os dados são sintéticos
(\fn{make\_sample\_data}), de modo que todo número abaixo é reproduzido
ao rodar \texttt{examples/tutorial.py}; com os seus dados, basta trocar o
Passo 1 e manter o resto.

O exemplo acompanha a região \texttt{L045}, que tem um número mediano de
especializações. Os rótulos de regiões e atividades são códigos
arbitrários.

\subsection{Passo 1: Da Tabela à Matriz}

A biblioteca parte de uma tabela longa: uma linha por região e
atividade, com um valor (emprego, exportações, patentes).
\fn{pivot\_to\_matrix} a transforma na matriz regiões $\times$
atividades que todas as funções recebem.

\begin{lstlisting}
df = ec.make_sample_data(n_locs=60, n_acts=40, seed=7)
print(df.head(3))
print(len(df), "rows |", df["loc"].nunique(), "regions |",
      df["act"].nunique(), "activities")

mat = ec.pivot_to_matrix(df, "loc", "act", "val")
print("matrix:", mat.shape)
\end{lstlisting}
\begin{lstlisting}[language={}, numbers=none, backgroundcolor=\color{white}, frame=leftline, framerule=2pt, rulecolor=\color{noteborder}, xleftmargin=1.5em]
    loc   act        val
0  L001  A001   2.450677
1  L002  A001  20.941576
2  L003  A001   3.358559
1450 rows | 60 regions | 40 activities
matrix: (60, 40)
\end{lstlisting}

Com o seu arquivo, a primeira linha vira
\texttt{df = pd.read\_csv("meus\_dados.csv")} e mudam os três nomes de
coluna. Vale conferir duas coisas antes de seguir, porque o pivot as
esconde: pares de região e atividade repetidos são somados, e pares
ausentes viram zero. Valores negativos ou faltantes são recusados pelas
funções, com um erro que nomeia o problema.

\subsection{Passo 2: Quem É Especializado em Quê}

Os valores brutos refletem sobretudo o tamanho: uma região grande
tem mais de tudo. A VCR (o quociente locacional, quando o dado é
emprego) divide a participação da atividade na região pela participação
dela na economia inteira, e a região conta como especializada onde
$\text{VCR} \geq 1$. A matriz binária $M$ registra essas
especializações; as somas das linhas são a diversidade de cada região e
as das colunas, a ubiquidade de cada atividade.

\begin{lstlisting}
rca = ec.rca(mat)
M = ec.mcp(mat)                    # 1 where RCA >= 1
diversity = ec.diversity(mat)      # specializations per region
ubiquity = ec.ubiquity(mat)        # specialized regions per activity

region = "L045"
print(rca.loc[region].nlargest(3).round(2))
print("specializations of", region, "=", int(diversity[region]))
print("median diversity =", diversity.median(),
      "| median ubiquity =", ubiquity.median())
\end{lstlisting}
\begin{lstlisting}[language={}, numbers=none, backgroundcolor=\color{white}, frame=leftline, framerule=2pt, rulecolor=\color{noteborder}, xleftmargin=1.5em]
act
A019    7.97
A009    3.62
A004    3.25
Name: L045, dtype: float64
specializations of L045 = 11
median diversity = 11.0 | median ubiquity = 17.5
\end{lstlisting}

A \texttt{L045} é especializada em 11 das 40
atividades, e a região mediana em 11. A sua
especialização mais forte tem VCR de 7{,}97: a
atividade pesa esse número de vezes mais na região do que na economia
como um todo. O limiar de 1 é o argumento \texttt{threshold} de toda
função que binariza, e \texttt{use\_rca=False} dispensa a VCR quando a
matriz já é um índice ou uma tabela 0/1.

\subsection{Passo 3: Complexidade, e Quanto Confiar Nela}

O ECI ordena as regiões pela complexidade daquilo em que são
especializadas, e o PCI ordena as atividades pela complexidade das
regiões que as têm. \fn{eci\_pci} devolve os dois; o padrão é o método
do autovetor, e \texttt{method=} troca para reflexos, Fitness ou GENEPY
na mesma chamada. \fn{eci\_diagnostics} diz em seguida quão firme é o
ranking nesta matriz em particular.

\begin{lstlisting}
eci, pci = ec.eci_pci(mat)
order = eci.sort_values(ascending=False)
print(order.round(2).iloc[[0, 1, -2, -1]])
print("ECI of", region, "=", round(eci[region], 2), "| rank",
      int(eci.rank(ascending=False)[region]), "of", len(eci))

diag = ec.eci_diagnostics(mat, n_boot=100, seed=1)
keys = ["gap", "corr_diversity", "rank_stability_mean", "sign_flips"]
print(diag[keys].round(3))
\end{lstlisting}
\begin{lstlisting}[language={}, numbers=none, backgroundcolor=\color{white}, frame=leftline, framerule=2pt, rulecolor=\color{noteborder}, xleftmargin=1.5em]
loc
L002    1.63
L005    1.54
L059   -1.94
L056   -2.38
Name: eci, dtype: float64
ECI of L045 = -0.22 | rank 35 of 60
gap                    0.115
corr_diversity         0.773
rank_stability_mean    0.954
sign_flips             0.000
Name: eci_diagnostics, dtype: float64
\end{lstlisting}

Os dois índices são escores z desta matriz: a \texttt{L045} tem ECI
de $-$0{,}22, posição 35
entre 60. Por serem relativos, o ECI de dois anos ou de duas
classificações não se compara em nível, só em posição; uma região ou
atividade devolvida como \texttt{NaN} ficou de fora por não ter nenhuma
especialização. Os diagnósticos aqui tranquilizam: o ECI tem correlação
de 0{,}77 com a diversidade e, quando um quinto das
atividades é retirado ao acaso, o ranking mantém correlação de postos de
0{,}95 com o original e nunca sai invertido. Em
matrizes pequenas ou esparsas esses números costumam ser bem piores. Se
a correlação com a diversidade for próxima de zero, o sinal do índice é
arbitrário, e convém conferi-lo com renda ou salários antes de ler ``ECI
alto'' como ``complexo''.

\subsection{Passo 4: O Que Está ao Alcance}

Duas atividades são próximas quando as mesmas regiões tendem a ser
especializadas nas duas. A partir da matriz de proximidade $\phi$, a
densidade de relatedness mede, para cada região e atividade, quanto do
que é relacionado à atividade a região já tem. É a medida básica da
biblioteca para a viabilidade de uma atividade numa região.
\fn{relatedness\_density\_external} guarda só as atividades que a
região não tem, que são as suas opções.

\begin{lstlisting}
phi = ec.product_proximity(mat)                  # proximity matrix
density = ec.relatedness_density(mat, phi=phi)   # 0 to 100
options = ec.relatedness_density_external(mat, phi=phi)

nearest = options.loc[region].nlargest(4)
print(nearest.round(1))

best = nearest.index[0]
why = ec.density_contributions(mat, region, best, phi=phi)
print(why.head(4).round(1))
print("sum of all contributions =", round(why.sum(), 1))
\end{lstlisting}
\begin{lstlisting}[language={}, numbers=none, backgroundcolor=\color{white}, frame=leftline, framerule=2pt, rulecolor=\color{noteborder}, xleftmargin=1.5em]
act
A017    38.1
A001    37.8
A005    35.6
A013    35.2
Name: L045, dtype: float64
act
A004    4.8
A014    4.0
A023    4.0
A003    3.8
Name: contribution, dtype: float64
sum of all contributions = 38.1
\end{lstlisting}

A atividade mais próxima que a \texttt{L045} não tem é a
\texttt{A017}, com densidade de 38{,}1: ponderando
pela proximidade, a região tem esse percentual do que é relacionado a
ela. \fn{density\_contributions} abre o número. Ela lista as atividades
da região que aproximam a \texttt{A017}, e as contribuições somam a
densidade. É a frase que entra na ficha de uma região: ``a
\texttt{A017} está ao alcance porque a região já tem
\texttt{A004}, \texttt{A014} e
\texttt{A023}''.

\subsection{Passo 5: O Que Vale a Pena Alcançar}

A viabilidade diz o que está ao alcance, não o que vale a pena
alcançar. O segundo eixo é uma escolha de quem analisa: complexidade,
salários, crescimento, emissões. \fn{desirability\_index} padroniza cada
indicador e tira a média ponderada, com os pesos declarados. A série de
salários abaixo é simulada e faz as vezes de uma fonte externa.

\begin{lstlisting}
rng = np.random.default_rng(7)
wage = 1000 * np.exp(0.3 * pci + rng.normal(0, 0.3, len(pci)))
indicators = pd.DataFrame({"pci": pci, "wage": wage})
weights = {"pci": 0.6, "wage": 0.4}

desirability = ec.desirability_index(indicators, weights)
robust = ec.weight_sensitivity(indicators, weights, top=10, seed=1)
print(robust.head(4).round(2))
\end{lstlisting}
\begin{lstlisting}[language={}, numbers=none, backgroundcolor=\color{white}, frame=leftline, framerule=2pt, rulecolor=\color{noteborder}, xleftmargin=1.5em]
      score  rank  rank_mean  rank_p05  rank_p95  top_share
act                                                        
A035   1.96     1       1.20       1.0       3.0        1.0
A038   1.61     2       3.42       2.0      10.0        1.0
A037   1.19     3       4.83       3.0       9.0        1.0
A040   1.17     4       5.20       4.0       7.0        1.0
\end{lstlisting}

Os pesos são um julgamento, e o ranking não deveria depender
deles. \fn{weight\_sensitivity} sorteia os pesos mil vezes em torno dos
escolhidos e registra onde cada atividade vai parar. A
\texttt{A035} é a primeira com os pesos escolhidos e fica
entre as dez melhores em 100\,\% dos
sorteios. Quando \texttt{rank\_p05} e \texttt{rank\_p95} ficam distantes,
a posição da atividade depende dos pesos, e o relatório deve dizer isso.
Um indicador que conta contra a atividade, como emissões, entra com peso
negativo.

\subsection{Passo 6: A Lista Curta}

Com os dois eixos, cada opção da região cai em um de quatro
quadrantes. Os cortes são as medianas das opções da própria região, de
modo que ``viável'' quer dizer viável em comparação com o que mais essa
região poderia fazer. \fn{opportunity\_alignment} resume o quadro num
número: a correlação entre os dois eixos nas opções da região.

\begin{lstlisting}
quadrants = ec.strategy_quadrants(options, desirability)
print(quadrants.loc[region].value_counts())

alignment = ec.opportunity_alignment(options, desirability)
print("alignment of", region, "=", round(alignment[region], 2))

shortlist = pd.DataFrame({
    "density": options.loc[region],
    "desirability": desirability,
    "quadrant": quadrants.loc[region],
}).dropna()
desirable = shortlist["quadrant"].isin(["short jump", "long bet"])
wanted = shortlist[desirable]
print(wanted.sort_values("density", ascending=False).head(4).round(2))
\end{lstlisting}
\begin{lstlisting}[language={}, numbers=none, backgroundcolor=\color{white}, frame=leftline, framerule=2pt, rulecolor=\color{noteborder}, xleftmargin=1.5em]
L045
low value     13
long bet      13
avoid          2
short jump     1
Name: count, dtype: int64
alignment of L045 = -0.87
      density  desirability    quadrant
act                                    
A024    30.24          0.91  short jump
A030    28.85          0.71    long bet
A026    28.40          0.52    long bet
A028    25.92          0.95    long bet
\end{lstlisting}

Das 29 atividades que a \texttt{L045} não tem, só
1 é salto curto (viável e desejável);
13 são apostas longas
(desejáveis, mas distantes) e
13 estão ao alcance,
mas valem pouco. O alinhamento de $-$0{,}87 diz o mesmo: nas
opções desta região, quanto mais próxima a atividade, menos desejável
ela tende a ser. É a situação comum de regiões de complexidade baixa ou
intermediária, e ela muda o que a lista curta precisa ser. Com tão
poucos saltos curtos, a lista útil é a metade desejável das opções
ordenada pela densidade: o único salto curto, a
\texttt{A024}, seguido das apostas longas menos
distantes, a \texttt{A030} e a
\texttt{A026}, com densidades entre
28{,}4 e 30{,}2. Numa
região
com alinhamento positivo, os saltos curtos encheriam a lista sozinhos.
Quantas apostas não relacionadas fazer, e em que ordem, é o que
\fn{optimize\_sequence} responde, e \fn{eci\_optimization} devolve o
conjunto de novas atividades que alcança um ECI-alvo com o menor esforço
(ver as seções de Difusão Estratégica e de Otimização de ECI).

\subsection{Passo 7: Lendo o Resultado por Listas de Política}

Políticas costumam ser escritas como listas de atividades: uma
missão, uma cadeia produtiva, uma taxonomia verde. As funções de lista
leem a mesma matriz por meio de uma lista. Abaixo, as quinze últimas
atividades fazem as vezes de uma missão e três grupos de atividades, dos
elos de uma cadeia.

\begin{lstlisting}
acts = list(mat.columns)
lists = {"mission": acts[25:]}
feasibility = density / 100            # shares, 0 to 1
tab = ec.list_complexity_potential(mat, pci, feasibility, lists)
print(tab.loc[region].round(2))

chain = {"inputs": acts[0:6], "components": acts[15:21],
         "assembly": acts[33:39]}
links = ec.missing_links(mat, chain, feasibility, cut=0.4)
print(links.loc[region])
\end{lstlisting}
\begin{lstlisting}[language={}, numbers=none, backgroundcolor=\color{white}, frame=leftline, framerule=2pt, rulecolor=\color{noteborder}, xleftmargin=1.5em]
         presence  potential  n_present  n_list
list                                           
mission       0.0       0.23          0      15
inputs        present
components    present
assembly      missing
Name: L045, dtype: object
\end{lstlisting}

A \texttt{L045} é especializada em
0 das
15 atividades da missão,
que aqui são as mais complexas. \texttt{presence} soma o PCI das
atividades da lista que ela tem, e \texttt{potential} é o PCI médio das
que faltam, cada uma ponderada pela sua densidade; com uma
lista de produtos verdes, são o índice e o potencial de complexidade
verde. Na cadeia, um elo é \texttt{present} quando a região é
especializada em pelo menos uma das suas atividades, \texttt{feasible}
quando alguma atividade do elo tem densidade de 40\,\% ou mais, e
\texttt{missing} caso contrário. \fn{transition\_risk} aplica a mesma
lógica a atividades que tendem a encolher.

\subsection{Passo 8: A Densidade Prevê a Entrada?}

Tudo até aqui supõe que a densidade diz para onde as regiões
vão. Com dois períodos isso pode ser testado: calculam-se os escores com
o primeiro período e verifica-se se eles ordenam as entradas observadas
até o segundo. Aqui o segundo período é simulado, com crescimento mais
rápido onde a densidade é alta; com dados reais, carregue o ano
posterior no lugar dele.

\begin{lstlisting}
size = np.outer(mat.sum(axis=1), mat.sum(axis=0)) / mat.values.sum()
mat_t1 = (mat * rng.lognormal(0, 0.3, mat.shape)
          + size * rng.lognormal(0, 1.0, mat.shape)
          * 3 * (density / 100) ** 2)

entered, at_risk = ec.entry_targets(rca, ec.rca(mat_t1))
print(int(at_risk.values.sum()), "cells at risk |",
      int(entered.sum().sum()), "entries")

scores = {"density": density,
          "relative": ec.relative_relatedness(mat, phi=phi),
          "persistence": rca}
val = ec.validate_entry(scores, entered, k=(5,),
                        baseline="persistence", n_boot=200, seed=1)
cols = ["auc", "auc_within", "precision@5"]
print(val["value"].unstack()[cols].round(3))
print(val.loc[("density", "auc")].round(3))
\end{lstlisting}
\begin{lstlisting}[language={}, numbers=none, backgroundcolor=\color{white}, frame=leftline, framerule=2pt, rulecolor=\color{noteborder}, xleftmargin=1.5em]
1400 cells at risk | 67 entries
metric         auc  auc_within  precision@5
score                                      
density      0.788       0.646        0.073
persistence  0.686       0.626        0.097
relative     0.639       0.646        0.073
value      0.788
lo         0.734
hi         0.839
diff       0.102
diff_lo    0.039
diff_hi    0.170
Name: (density, auc), dtype: float64
\end{lstlisting}

Uma célula está em risco quando a VCR é menor que 0,5 no início,
e entra quando a VCR chega a 1; o intervalo evita contar células que só
oscilam em torno de 1. De 1400 células em risco,
67 entraram. A AUC é a probabilidade de uma
entrada ficar acima de uma não entrada no ranking:
0{,}79 para a densidade e
0{,}69 para a persistência, que é
a VCR inicial. \texttt{auc\_within} calcula o mesmo dentro de cada região
e tira a média. É o número que importa para escolher entre as opções de
uma região, e é menor porque parte da AUC agregada vem de distinguir
regiões entre si: dentro das regiões, a densidade chega a
0{,}65 e a persistência a
0{,}63, uma distância bem
menor. \texttt{precision@5} é a parcela das cinco opções mais
bem colocadas de cada região que entraram. O último bloco compara a
densidade com a persistência: ganho de
0{,}10 na AUC, com intervalo de 95\,\%
de 0{,}04 a
0{,}17, obtido reamostrando regiões.
Para o teste ser honesto, todo escore tem de usar só dados do primeiro
período, inclusive a matriz de proximidade.

\subsection{Passo 9: Um Modelo de Previsão e o Emprego Esperado (Opcional)}

Com várias medidas de viabilidade, um modelo pode combiná-las.
\fn{build\_features} empilha as camadas numa tabela de células e
acrescenta controles; \fn{predict\_entry} faz validação cruzada por
região, de modo que a probabilidade de cada célula vem de um modelo que
nunca viu a sua região. Este passo exige o scikit-learn
(\texttt{pip install "econcomplex[ml]"}).

\begin{lstlisting}
layers = {"density": density, "implied": ec.implied_rca(mat)}
X = ec.build_features(mat, layers, at_risk=at_risk)
scores["model"] = ec.predict_entry(X, entered, model="logit", seed=1)
out = ec.validate_entry(scores, entered, k=(5,))["value"].unstack()
print(out.loc[["density", "model"], ["auc", "auc_within"]].round(3))

importance = ec.explain_entry(X, entered, model="logit", seed=1)
print(importance.head(3).round(3))

# next window: fit on t0 -> t1, predict from what is observed at t1
phi_t1 = ec.product_proximity(mat_t1)
layers_t1 = {"density": ec.relatedness_density(mat_t1, phi=phi_t1),
             "implied": ec.implied_rca(mat_t1)}
X_t1 = ec.build_features(mat_t1, layers_t1,
                         at_risk=ec.rca(mat_t1) < 0.5)
prob = ec.predict_entry(X, entered, model="logit",
                        new_features=X_t1, calibrate=True)
jobs = ec.potential_employment(prob, mat_t1)
print(jobs.groupby("region")["expected"].sum().nlargest(3).round(1))
\end{lstlisting}
\begin{lstlisting}[language={}, numbers=none, backgroundcolor=\color{white}, frame=leftline, framerule=2pt, rulecolor=\color{noteborder}, xleftmargin=1.5em]
metric     auc  auc_within
score                     
density  0.788       0.646
model    0.792       0.643
                 importance    std  p_value  significant
density               0.273  0.099    0.020         True
log_region_size       0.040  0.037    0.137        False
rca                   0.023  0.021    0.137        False
region
L019    540.4
L016    530.3
L017    519.1
Name: expected, dtype: float64
\end{lstlisting}

O modelo entra em \fn{validate\_entry} como mais um escore. Aqui
ele empata com a densidade pura (AUC de 0{,}792
contra 0{,}788, e
0{,}643 contra
0{,}646 dentro das regiões), que é o que
se espera de dados simulados só a partir da densidade. Um modelo só compensa a maquinaria a mais
quando ganha essa comparação fora da amostra. \fn{explain\_entry}
embaralha um atributo de cada vez nas regiões fora do treino e informa a
queda da AUC: a densidade sustenta a previsão
(0{,}27) e nenhum outro atributo se
distingue estatisticamente de zero. Para a janela seguinte, a mesma função recebe
\texttt{new\_features}: o modelo é ajustado na primeira janela e
aplicado ao que se observa no fim dela. \fn{potential\_employment}
multiplica então cada probabilidade calibrada pelo tamanho que a
atividade teria na região com quociente locacional 1, e a soma por
região é o emprego esperado em novas especializações.

\subsection{Antes de Usar os Resultados}

Alguns pontos respondem pela maior parte dos erros no uso aplicado
desses indicadores.

\begin{itemize}[leftmargin=1.5em]
  \item ECI e PCI são relativos a uma matriz. Para acompanhar uma região
    no tempo, compare posições, ou projete um PCI fixo em cada ano com
    \fn{subnational\_eci}.
  \item \texttt{NaN} quer dizer que a unidade ficou de fora (sem
    especialização, ou não é uma opção); não é um zero.
  \item Em dados esparsos, em que muitas regiões têm uma ou duas
    especializações, calcule a complexidade no núcleo conectado, com
    \texttt{dmin=2, umin=2}.
  \item A densidade vem em percentual (0 a 100) por padrão e como
    parcela com \texttt{percent=False}. Cortes como o de
    \fn{missing\_links} são lidos na escala da matriz que você passa.
  \item Vetores externos (PIB, população, salários, PCI de outra fonte)
    são casados com a matriz pelo rótulo, em qualquer ordem; um rótulo
    ausente gera erro, em vez de ser preenchido.
  \item Em qualquer teste de previsão, a matriz de proximidade e todos
    os escores usam só o primeiro período.
  \item Declare os pesos do índice de desejabilidade e informe como o
    ranking se move quando eles mudam.
\end{itemize}
